\section{Startup and unbounded binary simulation}
\label{sec:runtime}

\subsection{Exact blank startup}

Direct execution of the released machine from the all-zero tape reaches
the complete expected dispatch frame after exactly 409,835 transitions.
The canonical control state is 22 and the head is at $-5$. The logical
primary registers are
\[
                     (3,0,0,0,3,0,0,0,0).
\]
The logical-to-physical register map is
$(6,2,0,4,3,7,8,5,1)$, giving physical primary values
$(0,0,0,3,0,0,3,0,0)$. There are 167 spare register blocks, each
encoding one. The zero-gap parameter is 181, or 117 beyond the
runtime checker's fixed 64-cell address allowance. Exactly 351 cells
contain one; their complete positions are in the startup receipt, and
every other tape cell is zero.

The builder computes this endpoint with a sparse set of one-cells.
An independent dense replay of the canonical table computes the same
entire tape support, head, and control state. Its finite allocated
window is only a computational implementation detail: leaving that
window causes the check to fail, and no such event occurred. The
infinite exterior therefore remains exactly zero. This finite startup
check is exhaustive for this particular initial segment. It is separate
from the unbounded theorem for subsequent instruction frames.

\subsection{What the universal runtime checker establishes}

Each primitive source instruction has a binary address. For each such
address, the checker starts with its literal address in the dispatch
frame, nine symbolic natural primary values, a symbolic extra zero gap,
a symbolic first spare value, an arbitrary finite remaining spare list,
and zero tape exteriors. It derives the exact successor frame or halt.
For a decrement instruction it considers the zero and positive cases
separately; the positive case represents its register as $n+1$ with
arbitrary natural $n$.

The proof engine uses three kinds of justified steps:
\begin{enumerate}
\item A literal transition from the supplied binary table.
\item Induction over a natural run length when the literal transition
      is a self-loop scanning in the direction of motion.
\item Structural induction over a finite list of unary blocks, after
      deriving its one-block schema from this exact table.
\end{enumerate}
The four list traversals are rederived, rather than accepted as
unverified macro instructions. Unary blocks have forms
$B(v)=0\,1^{v+1}$ and $C(v)=1\,0\,1^v$. The tape normalizer uses
their proved concatenation identities, including
$C(A)1=1B(A)$ for an arbitrary finite block list $A$ and its reversal.
Opaque tape contexts cannot be inspected without an applicable proved
rule. The checker rejects unknown reads and unclosed symbolic cases.

The primary candidate passed all 6,167 guarded cases. It used
1,488,390 symbolic rule applications, covered all 444 literal rows,
and established a positive affine concrete duration for every
instruction, with minimum constant term 13. Thus no simulated source
step can disappear into a zero-time loop or an unproved divergent
microprogram. The symbolic iteration cap limits the proof search the
checker can complete; it never accepts a case just because the cap
was reached. Natural register values and concrete running times are
not bounded by that cap.

\begin{center}
\begin{tabular}{lrr}
\toprule
Check & Seed plus $G$ & Conditional K plus $G$\\
\midrule
Working states & 222 & 222\\
Literal rows & 444 & 444\\
Guarded instruction cases & 6,167 & 5,956\\
Symbolic rule applications & 1,488,390 & 1,443,614\\
Rows covered & 444 & 444\\
Minimum concrete duration & 13 & 13\\
Blank startup transitions & 409,835 & 409,835\\
\bottomrule
\end{tabular}
\end{center}

All-row coverage is useful evidence that every literal row has been
checked in some supported instruction context. It is not by itself a
proof of the global behavior. That proof uses the frame invariants,
exact endpoints, and positive progress at every simulated instruction.
The executable checker is an ordinary mathematical verification program;
its result is not a kernel-checked theorem in Rocq or Lean.

\section{From the actual enumerator to consistency}
\label{sec:chain}

The finite transition folds and runtime certificates establish the
implementation layers. To obtain the advertised consistency equivalence
one must also know exactly what proof search the unchanged register
program performs and why its nonzero initialization still searches all
relevant proofs. This section records those inherited interfaces and
their application to the changed source.

\subsection{The changed row does not disturb proof extraction}

The proof commands act on a persistent tuple $(t,w,d,h)$ consisting of
the current formula, an encoded stack, an additive parameter counter,
and a cached formula. An axiom command chooses its parameters through
$a=d+x$, $b=h+y$, and its explicit $c,u$ fields, resets $d,h$ to zero,
returns the gated bundle, and pushes the eight logical rows together
with the generalized old cache and the old current formula. The exact
ordering is retained, including the replacement of row $\rho_7$ by $G$
at the same position.

The soundness invariant says that the current formula, cache, and
stack contents are theorems of the specified numerical theory. The
truth seed is $0$. Axiom commands preserve theoremhood because their
outputs are exactly the newly specified axiom instances; MP, DROP,
NIP, and theorem generalization keep their old meanings. A halting
NIP has cached code $1$, giving a refutation. All logical rows and
enabled bundles have root constructors that exclude accidental code
$1$; the disabled bundle is $0$.

For the reverse implication the inherited proof compiler gives, for
each finite derivation of formula $X$, a command block with the contract
\[
 (T,w,0,0)\ \longmapsto\ \text{halt or }(X,w,0,0)
 \quad\text{for every }T,w.
\]
Its leaf-selection blocks use the positions of the eight logical rows,
the unchanged disabled gate, and the fact that skipped rows are not
$1$. These properties still hold. An arbitrary $G(a,b,c)$ instance
can be selected by the axiom fields after the reset block makes $d=h=0$;
the gate on the separately emitted bundle does not restrict this logical
row. The bundle leaf now returns exactly~\eqref{eq:finalbundle}.
The GEN and MP block proofs are unchanged. The corresponding pinned
formal interfaces are \code{Enum230.v} and \code{L1_Comp230.v}; their
precise ordinary adaptation is retained in the K224 packet.

\subsection{A cofinal family handles the nonzero search cutoff}

It would be insufficient to say that the enumerator eventually checks
every small proof: its initialized next-proof counter is three and it
retains state between candidates. The inherited construction provides
an explicit way to obtain sufficiently large halting candidates.

Let $\mathrm{Build}_F$ be the block for a derivation of $F=1$. For each
$U\in\mathbb N$, use the resetting command
\[
 \mathrm{START}_U=\mathrm{AX}(0,1,1,U)
\]
and follow it by $\mathrm{Build}_F$ and
$[\mathrm{DROP},\mathrm{NIP},\mathrm{NIP}]$.
The reset command makes $d=h=0$ regardless of the retained state.
It may change the stack and current formula, but the block contract
holds for every such new pair. If the block has not already halted,
the final three commands place $F$ in the cache and halt.

If $R$ is the fixed encoded tail after the reset prefix, the complete
candidate number is exactly
\[
 N_U=P\bigl(3,P(1,P(1,P(U,R)))\bigr).
\]
Since $P$ is at least either argument, $N_U\geq U$. Thus this is a
cofinal family of halting candidate numbers, and $U\geq3$ suffices
for the actual cutoff. The drain proof makes the positive residual
command-list code decrease within a candidate; at a boundary the
next-proof counter increases. Hence an appropriate $N_U$ is eventually
started unless an earlier candidate already halts. The reset deals
with any earlier retained state.

This is a finite, explicit completeness argument for the actual seeded
enumerator. It does not assume that a fixed finite command word has
arbitrarily large encodings without exhibiting them.

\subsection{The complete proof chain and its trust boundary}

The following equivalences compose:
\begin{align*}
 M_{\mathrm{sG}}\text{ halts from blank}
 &\longleftrightarrow
    \text{the delivered register program halts from its seed},\\
 &\longleftrightarrow T_{\mathrm{sG}}\text{ derives native code }1,\\
 &\longleftrightarrow T_{\mathrm s}\text{ derives native code }1,\\
 &\longleftrightarrow \mathrm{ZFC}\text{ has a finite refutation}.
\end{align*}
The first line uses complete startup and the unbounded runtime
simulation of the independently identified source. The second is
the proof-enumerator argument above. The third is
Theorem~\ref{thm:deductive}. The fourth is the inherited seed/numerical
calibration detailed in Sections~\ref{sec:seed} and~\ref{sec:native}.
Negating these finite-existence statements gives
Theorem~\ref{thm:main}.

Every new operation is finite syntax manipulation, induction over a
finite derivation or list, induction over a natural block length, or
an exact finite execution. The stated inherited interpretations and
proof transformations are also effective on finite proofs. These
arguments admit the ordinary arithmetization in ZFC required by the
user's target. This is distinct from claiming a completed internal PA
formalization, or from claiming that a Python run proves the consistency
of ZFC.

\begin{center}
\begin{tabular}{>{\raggedright\arraybackslash}p{0.32\linewidth}>{\raggedright\arraybackslash}p{0.58\linewidth}}
\toprule
Layer & Evidence and scope\\
\midrule
New equality/bundle comparison & Uniform proof and newly checked closed
proof DAG for the exact seed order.\\
New source & Pinned inputs, symbolic expression identities, arbitrary-register
macro bisimulations, full instruction-word and guarded-CFG checks.\\
New table & Deterministic construction, exact row census and relabeling,
two independently certified all-tape folds.\\
Runtime and startup & Fresh all-natural instruction proofs on the literal
table; sparse and independent dense blank execution.\\
Inherited logical strength & K230/K224/seed223 finite-proof interpretations,
exceptional-code bridge, equality interpretation, and K0-to-ZFC calibration.\\
Proof-assistant status & No newly compiled end-to-end Rocq or Lean endpoint
for either 222-state table.\\
\bottomrule
\end{tabular}
\end{center}

The inherited dependencies are real mathematical premises of this
construction. The new checks do not erase them, and retaining old
formal sources or receipts is not the same as replaying every old
kernel build. The reproduction instructions distinguish the new
replay from an optional predecessor replay.

\section{The second 222-state construction}
\label{sec:alternative}

A parallel construction begins with the published conditional-K223
source and applies the same $G$ absorption. Its old core has the
conditional root
\[
      \varphi(x,a)\to\exists C\,(R(x,C)\land H_\varphi(C)).
\]
The helper is removed from its bundle and the unconditional equality
row is replaced by~\eqref{eq:g}. Exact symbolic macro execution and
both complete guarded-CFG resolvers bind this new source to its stated
formula. It has 4,581 primitive occurrences, 209 distinct subtrees,
and maximum zero run six.

Compilation initially gives 224 states. Two independently checked
one-step folds identify historical state 135 with 118, then state 116
with 159. The same kind of old scanned-zero invariant supports the
first fold; the second is an identical-row quotient. The final table
has 222 working states. Its K-labelled form has SHA-256
\begin{center}\small\ttfamily
7f0d5c464b738046fceacb8219d21e17b66e5820160ef2c0e984a4281c20faf5.
\end{center}
The archived alternative includes a canonical labeling, source, maps,
checkers, and its own runtime and startup receipts. Its canonical file is
\path{portable/alternative/results/candidate222.tm}, with SHA-256
\begin{center}\small\ttfamily
60e21a040e58e2a2840b0f3d58ad3ee0e6ef4d0d5cdf1372d5a69355f1198ca1.
\end{center}

Its 5,956 guarded runtime cases passed with 1,443,614 symbolic rule
applications and all 444 rows covered. Sparse and dense startup
reached the same intended frame after 409,835 transitions. The $G$
interderivation is uniform in the set-theoretic core, so it transfers
the published conditional candidate's ordinary calibration unchanged.
That calibration includes an additional two-marker hereditary-domain
interpretation relating conditional and unconditional K. It remains
an explicitly inherited ordinary proof, with no new kernel discharge
in this packet.

The seed-plus-$G$ version is the primary recommendation because its
set-theoretic comparison uses the existing seed-to-K0 route directly.
The conditional variant is useful independent evidence that the
equality saving composes with two structurally different K encodings.
They have equal state counts and distinct sources and transition tables.
